\section{Introduction}
\label{sec:intro}

A column store of the MergeTree family holds a table as a set of immutable
parts, each sorted by the table's ordering key and cut into granules of a few
thousand rows, with a sparse primary index and a set of skipping indexes
written inside the part~\cite{schulze2024clickhouse,clickhousedocs}. Retrieval
that ranks rows by similarity to a query fits that layout awkwardly. The text
index answers which rows contain a token, not which rows rank highest for a
weighted query. The vector index is built per part and searched per part, so a
similarity query visits every part that survives the predicate; the same
arrangement holds in Lucene and in Milvus, where a segment is the unit an
approximate index is built over~\cite{lucene,wang2021milvus}. As a table grows,
the number of parts grows, and the cost of a similarity query grows with it
unless something excludes parts before they are searched.

This paper describes a design for that gap and states how we intend to measure
it. Two structures carry it. The first is a weighted inverted file over a mixed
sparse basis: each basis element --- a word, a linked entity, a code symbol, a
coordinate of a learned sparse dictionary, a relation, a structural motif ---
owns a posting list of (row, weight) pairs, and a query written over the same
basis is scored only over the posting lists of the elements it activates, with
top-$k$ evaluation under block-max pruning. The second is a routing structure
built per part: regions with a centroid, a covering radius, the maximum weight
each basis element attains in the region, and counting statistics, aggregated
into a table-level structure that the planner consults to drop parts whose best
possible score cannot reach the running $k$-th best, before any per-part index
is opened.

\paragraph{What is claimed.}
Expected retrieval cost should scale with the relevant neighbourhood rather
than with corpus size: at a fixed query distribution and a fixed recall target,
the parts, granules and bytes a query touches should grow more slowly than the
corpus does, and should stay within a bounded factor of the granules that hold
the answer. That is an empirical claim about a distribution of queries, and
\S\ref{sec:method} is the protocol that would support or refute it.

\paragraph{What is not claimed.}
No worst-case guarantee is available here and none is asserted. A query whose
relevant set is the corpus --- $k$ comparable to $N$, or a request that must
enumerate every match of a basis element every row carries --- costs time
linear in the corpus, under this design and under any other. A query whose
region bounds all exceed the $k$-th best score degrades the routing to a visit
of every part, which is the cost of the baseline plus the cost of having
consulted the summaries. High intrinsic dimension makes that degradation the
common case rather than the rare one~\cite{weber1998quantitative,beyer1999nn},
and graph indexes have their own hard instances~\cite{indyk2023worst}. The
design is an expected-cost design; the protocol reports the tail as well as the
mean so that a claim about the mean cannot hide behind it.

\paragraph{Status of the claims.}
Sections~\ref{sec:setting}--\ref{sec:routing} describe a design and label each
component \prior{} or \proposed{}; nothing in them is a measurement.
Section~\ref{sec:method} specifies a measurement that has not been made.
Section~\ref{sec:results} holds the table it will fill, and every cell in it
prints \ph{} until the harness writes \texttt{results-data.tex}.
Section~\ref{sec:related} places the design against published work, and
Section~\ref{sec:limits} lists what would make the whole approach the wrong
one.

\paragraph{Contributions.}
\begin{itemize}[leftmargin=*]
\item A retrieval index over a heterogeneous sparse basis, laid out inside the
  parts and granules of a merge-based column store, with block maxima aligned
  to granules so that the store's existing skipping unit is also the pruning
  unit (\S\ref{sec:index}).
\item A part-level routing structure and the planner rule that reads it, built
  from bounds that metric indexes and dynamic pruning have used separately, and
  maintained through the merge lifecycle that produces the parts
  (\S\ref{sec:routing}).
\item A measurement protocol for expected cost, which holds recall fixed while
  corpus size varies, fits an exponent to each cost counter, reports bytes
  fetched per correct answer, and separates the queries for which the design is
  expected to fail (\S\ref{sec:method}).
\item A statement of which parts of the design are already published work, with
  the references, and which are not (\S\ref{sec:related}, Table~\ref{tab:status}).
\end{itemize}
